\section*{結言}
この日のXでは、巨大なオープンウェイトモデルの発表と、学習データ・パイプラインの公開が並んだ。Kimi K3.1 Proの性能主張とK2 Horizonのデータ公開は、モデルの能力だけでなく、研究成果をどの範囲まで再利用できるかという競争を示している。推論基盤llmdの更新も、モデルの実用化を支える動きとして位置づけられる。

一方、OpenAIの研究環境分離とAnthropicの逸脱行動分析は、安全性を運用と評価の両面から扱う必要性を示した。Geminiの科学応用は成果の可能性を示すが、個別の主張には独立検証が残る。Canva、Grok Bot、Mubitに関する話題も公式資料が未確認の部分を明示して読み、今後の仕様公開や実利用の報告を追いたい。
